\begin{lemma}[Homogeneous triples]
For every infinite $X\subseteq\mathbb N$ and every two-coloring of the
three-element subsets of $X$, there are an infinite $Y\subseteq X$ and a
color $b$ such that every three-element subset of $Y$ has color $b$.
\end{lemma}
\begin{proof}
Take the front $F=\mathsf{CardinalityFront}(3,X)$ from
\noderef{CardinalityFront}; it is a front by
\noderef{CardinalityFrontIsFront}.  Let
$S=\{s\in F\mid c(s)=1\}$.  Apply the Nash--Williams theorem
\noderef{NashWilliams} to $F$ and $S$.  It gives a front $F'$ on an
infinite $Y$, with $F'\subseteq F$, such that either
$F'\subseteq S$ or $F'\cap S=\emptyset$.

Since $F'$ is nontrivial (all its members have cardinality three), its
base is the union of its members.  Every member of $F'$ belongs to $F$,
and every member of $F$ is contained in $X$, so this also gives
$Y\subseteq X$.

Because every member of $F'$ has cardinality three, front density forces
$F'$ to contain the first three elements of every infinite branch through
$Y$: for a given three-element $s\subseteq Y$, take the infinite set
whose first three elements are exactly $s$, and apply density of $F'$;
the member supplied has cardinality three and hence is $s$.  Thus every
three-element subset of $Y$ lies in $F'$.  In the first homogeneous case
all have color $1$, and in the second all have color $0$.  Choosing the
corresponding Boolean value gives the conclusion.
\end{proof}
